% ECONOMICS_PROBLEM: {"id":"D82-46","title":"Robust dynamic mechanisms","jel_code":"D82","source_id":"OP-0290"}
% FIRST_POSTED: 2026-10-10
% ATTEMPT_STATUS: solved
% ENTRY_KIND: research_attempt
\documentclass[11pt,a4paper]{article}
\usepackage[left=20mm,right=20mm,top=18mm,bottom=18mm]{geometry}
\usepackage[T1]{fontenc}
\usepackage{lmodern,microtype}
\usepackage{amsmath,amssymb,amsthm,mathtools}
\usepackage{xurl,hyperref}
\hypersetup{hidelinks,hypertexnames=false,pdftitle={D82-46: Exact robust dynamic mechanism design}}
\setlength{\parindent}{0pt}
\setlength{\parskip}{3pt}
\setlength{\emergencystretch}{2em}
\allowdisplaybreaks
\raggedbottom
\widowpenalty=10000
\clubpenalty=10000
\displaywidowpenalty=10000
\newtheorem{theorem}{Theorem}
\newtheorem{lemma}[theorem]{Lemma}
\newtheorem{proposition}[theorem]{Proposition}
\newcommand{\R}{\mathbb R}
\title{D82-46: Exact robust dynamic mechanism design\\with finite types and continuous payments}
\author{Ji Ho Bae}
\date{10 October 2026}
\begin{document}
\maketitle
\begin{abstract}
Finite types do not make a dynamic mechanism finite: publicly observed
real-valued payments can support infinitely many continuation histories.
We prove that every committed robust mechanism has exactly the revenue
function of a bounded finite public-outcome tree, retaining all laws,
allocation-controlled transitions and adapted reporting-and-exit deviations.
An explicit polynomial recursion then characterizes robust truthfulness,
participation and the exact worst-case revenue ratio. For the catalogue's
finite or polynomially specified marginal, transition and moment ambiguity
classes, real quantifier elimination computes the ratio and implementing
mechanisms, with a separate test of attainment. A distribution-independent
positive guarantee and a two-period known-marginal example show respectively
when positive revenue survives and when belief ambiguity removes gains from
intertemporal screening.
\end{abstract}

\section*{Target, method, and claim status}
\textbf{Model:} GPT Codex. \textbf{Contributor:} Ji Ho Bae.
The exact serving revision was not exposed. The construction and hand proof
were produced through parallel derivation and independent adversarial review.
No human-written mathematical edits or proof-assistant verification are
represented. Expert review and catalogue acceptance remain separate.

The complete original generated output is preserved at the immutable
\href{https://github.com/jbaelaw/robust-dynamic-mechanism-characterization/blob/40214a615d83117a2b598613cc314aaba833ca10/paper/Proof.tex}{manuscript TeX}
and \href{https://github.com/jbaelaw/robust-dynamic-mechanism-characterization/blob/40214a615d83117a2b598613cc314aaba833ca10/paper/Proof.pdf}{manuscript PDF}.

The target is the finite-horizon editorial specification D82-46, with finite
types and allocations, bounded valuations and payments, public report histories,
and an explicitly declared continuation and participation convention.
The answer retains continuous payment signals and does not replace them by
mean payments. The finite support bound in Section~2 is the essential reduction;
the later real-algebraic procedure has no asserted polynomial-time bound.
The structural characterization holds for arbitrary ambiguity sets. Effective
computation requires the stated finite semialgebraic representation and exact
rational or real-algebraic input. Section~5 gives sufficient restrictions for a
positive guarantee; Section~6 proves a strict loss with known marginals.
These examples accompany the general characterization rather than substituting
for its original quantifiers.

Bergemann and V\"alim\"aki's June 2018 revision, Section~8, printed
pp.~48--50, motivates weaker common-prior assumptions and continuation
participation~\cite{BV}. The catalogue's exact minimax ratio and finite model
are editorial specifications, not results attributed to that survey.
The proof below is self-contained apart from standard real quantifier
elimination~\cite{Collins}; no historical priority claim is made.

\section{Finite primitives and a precise continuation convention}
\label{sec:primitives}

Fix the following finite primitives, independently of the mechanism.
There are $n$ agents, dates $1,\ldots,T$, a finite joint type set
$\Theta=\prod_i\Theta_i$ of cardinality $q$, and a finite nonempty
allocation set $X$. Values satisfy $0\leq v_i(x,\theta_i)\leq1$.
The discount factor is $0<\beta\leq1$ and transfers lie in $[-B,B]$,
where $0\leq B<\infty$. Feasibility may alternatively restrict the
allocation paths to a declared finite set with known feasible prefixes;
all tree shapes below are then filtered by those prefixes.

A law $P$ is an initial joint distribution $\pi_P$ and a collection of
controlled kernels
\[
 K_{P,t}(\theta^{t+1}\mid\theta^{1:t},x^{1:t}),
 \qquad 1\leq t<T.
 \tag{F1}
\]
Every coordinate of $\pi_P$ and these kernels is strictly positive and
the appropriate rows sum to one. Thus arbitrary correlations, memory,
and allocation effects are allowed. Nature does not condition on a
payment signal except through the declared finite allocation history.
Write $z=(z_1,\ldots,z_d)$ for the finite vector of law coordinates.
An ambiguity class $A\subset\mathbb R^d$ is fixed before choosing a
mechanism and satisfies the displayed probability restrictions.
For the effective algebraic conclusions, $A$ is specified by a finite
Boolean combination of polynomial equalities and inequalities.
The finite representation and incentive characterization below hold
for an arbitrary nonempty $A$; a terminating algebraic algorithm is
claimed only for the stated semialgebraic input class. For an arbitrary
$A$ the result is an exact finite-variable, possibly semi-infinite
constraint system, with no effective-algorithm assertion without an
effective presentation of that set.

At the beginning of date $t$, agent $i$ knows its own type history and
the public report, allocation, and transfer history from earlier dates.
Current reports are simultaneous; their joint profile becomes public
before the allocation and transfer lottery is realized. A mechanism is
a committed Borel stochastic map from these histories and current
reports into $X\times[-B,B]^n$. Its public history can include arbitrary
past real-valued transfers. Private correlated mechanism randomization
is allowed: it can be represented by its conditional stochastic kernels
given its observable history, without exposing its future random seed.

Exit is available before the current report, gives zero future utility,
and leaves earlier transfers sunk. After any public history, opponents
are regarded as truthful. Hence agent $i$'s true private history together
with the opponents' past public reports determines the full past true
type history. Its distribution for the current opponents' types is the
kernel in \textup{(F1)}, conditioned on its own current type. This is also
the prescribed belief at off-path report histories: past physical
allocations are treated as fixed controls, and the same strictly positive
type kernels are used. Thus no division by the probability of an
off-path report or a singleton real-valued transfer signal is performed.
This convention agrees with Bayes' rule whenever that history has
positive probability under truthful opponents, including after one's
own earlier false reports. It gives a defined continuation problem at
every private history.

To check that agreement, fix the full true past reconstructed by agent
$i$ and a public allocation, payment, and report history. The mechanism
likelihood of that history depends on its public report input and its
own randomization, and not on an additional hidden true type. The
conditional type likelihood therefore factors as the controlled nature
likelihood times a common mechanism likelihood. The latter cancels.
This argument also applies to regular conditional kernels for continuous
payments, outside a null set; the displayed pointwise prescription
specifies the remaining histories. Arbitrary payment-dependent beliefs
or payment-controlled nature kernels would be different primitives.

Truthfulness means optimality of truthful continuation against every
adapted reporting-and-exit strategy at all those private histories.
Participation means its value is nonnegative there. The same information,
physical, transfer, exit, and off-path conventions define the known-law
benchmark; only the set of laws in its incentive inequalities is reduced
to the singleton $\{P\}$.

\section{An exact finite public-branching reduction}
\label{sec:compression}

The transfer interval is not discretized. The following support bound
keeps the entire continuous-transfer mechanism class.
Set
\[
 N_t=\binom{d+T-t}{T-t},\qquad
 D_t=(n+1)q^tN_t,\qquad b_t=D_t+1.
 \tag{F2}
\]
These are upper bounds, not claims of minimal size.

\begin{lemma}[Finite barycentres without a closure gap]
\label{lem:barycentre}
Let $Y$ be a bounded measurable random vector with values in a subset
$S$ of $\mathbb R^m$. Its expectation is a convex combination of at
most $m+1$ points of $S$.
\end{lemma}
\begin{proof}
Write $C=\operatorname{conv}(S)$. Separation first implies
$\mathbb EY\in\overline C$: a separating affine functional otherwise
contradicts taking its expectation. If $\mathbb EY$ belongs to the
relative interior of $\overline C$, it belongs to $C$. Indeed the
relative interior of a convex set and that of its closure agree, as
can be seen by approximating the vertices of a small simplex surrounding
the point by points of $C$.

Otherwise choose a supporting affine hyperplane of $\overline C$ at
$\mathbb EY$ in its affine hull. Its defining affine functional is
nonnegative on $S$ and has expectation zero. It is therefore zero
almost surely. Apply the same argument in this hyperplane to the
almost-sure range of $Y$. The affine dimension decreases, so after
finitely many repetitions the expectation lies in the convex hull of
the corresponding actual range. Finally, any convex combination with
more than $m+1$ points can be shortened: their augmented vectors
$(s,1)$ are linearly dependent, and one moves along that dependence
until a coefficient becomes zero. Iteration proves the bound.
All selected points are actual values, rather than limiting points
that might have different observable payments.
\end{proof}

\begin{lemma}[Exact public-outcome compression]\label{lem:compression}
For every admissible mechanism $M$, there is an admissible mechanism
$\widehat M$ with at most $b_t$ observable allocation--payment outcomes
after each public history and report profile at date $t$, such that:
\begin{enumerate}
\item truthful expected utility at every retained public history is
unchanged for every full actual type history and every controlled law;
\item every one-step reporting deviation has unchanged continuation value;
\item truthful revenue is unchanged for every controlled law.
\end{enumerate}
In particular $M$ and $\widehat M$ have identical revenue functions on
$A$, and robust truthfulness and participation are preserved.
The continuation kernel at a selected outcome is retained, rather than
its future random seed being revealed.
\end{lemma}
\begin{proof}
Fix a public history $h$, current report profile $r$, and a realized
current outcome $(x,p)$. Keep the original continuation at that exact
observable history. For every possible full actual type history
$\eta=\theta^{1:t}$, consider each agent's current-plus-future utility
when all future reports are truthful, and also total current-plus-future
transfers. Regard these quantities as functions of the controlled law
coordinates $z$.

Each function is a polynomial of degree at most $T-t$. To see this
directly, sum over the finitely many future type and allocation paths.
For a fixed such path, the nature likelihood is a product of at most
$T-t$ kernel coordinates. Its coefficient is the mechanism probability
of the allocation path times the conditional expected discounted value
or transfer, integrated over its possibly continuous payment signals.
The mechanism coefficients do not depend on $z$: its report input is
fixed by the specified true future type path, and past realized
allocations have already fixed the arguments of nature's kernels.
There are finitely many monomials and all coefficients are bounded.
Equivalently the same conclusion follows by backward expansion of
the stochastic kernels and bounded Fubini integration.

Collect all monomial coefficients of these $n+1$ functions for all
$q^t$ possible $\eta$ into one vector. It lies in
$\mathbb R^{D_t}$, with harmless zero coordinates if some monomials
are absent. Its dependence on the current outcome is measurable and
bounded. Lemma~\ref{lem:barycentre} replaces the current outcome
distribution by a distribution on at most $D_t+1$ actual observable
outcomes with exactly the same coefficient vector. At each selected
outcome use its original continuation kernel. Repeated selected
outcomes are combined; there is then a unique continuation for every
observable selected $(x,p)$.

Do this at the root separately for every report profile, then at the
finitely many selected children, and continue to date $T$. Subsequent
compression of a selected child preserves all its continuation
polynomials, so it leaves each already preserved parent polynomial
unchanged. This top-down procedure uses only finitely many histories
at each of the finitely many dates.

Why do these polynomial identities preserve incentives? At a pre-report
private history the actual past joint type path is known under truthful
opponents. For a fixed own current type, the current opponents' types
have conditional probabilities equal to the relevant positive initial
or transition coordinates divided by their sum. Those coordinates do
not depend on the past real-valued payment signals. A one-step false
report followed by truth therefore has a value that is a finite weighted
sum of exactly the preserved continuation polynomials. Histories after
one's earlier false reports are included because every full actual
history $\eta$, not merely those matching past reports, was preserved.
Truthful values and immediate-exit values are preserved in the same way.

For completeness, one-step truthfulness and nonnegative truthful value
at every history imply optimality against every adapted reporting-and-exit
strategy. At the final date this is the one-step comparison. Inductively,
after any current report and observable outcome, replacing a deviating
future strategy by truth weakly improves its conditional payoff; compare
the remaining current report with truth. Immediate exit yields zero.
Bounded payoffs justify all conditional expectations. Randomized
deviations are mixtures of these comparisons. This proves the induction.

Public histories outside the finite selected outcome tree have zero
probability even after arbitrary reports. Complete the mechanism there
by any fixed feasible report-independent allocation rule and zero
transfers for the remaining dates. Values are nonnegative, so this is
truthful with voluntary continuation. If feasibility depends on the
physical prefix, use a declared feasible report-independent extension
of that prefix; infeasible histories are not private histories of the
allocation problem. This completes the off-path prescription and the
proof.
\end{proof}

\section{Explicit finite polynomial incentive constraints}
\label{sec:polynomial}

Enumerate finite rooted tree shapes having, after every report profile
at date $t$, between $1$ and $b_t$ outcome edges. Each edge has a fixed
allocation label $x\in X$, a probability $w>0$, and a transfer vector
$p\in[-B,B]^n$; sibling probabilities sum to one. Positive-probability
siblings must have different observable pairs $(x,p)$. For siblings
with the same fixed allocation label, this is the polynomial condition
$\sum_i(p_i-p'_i)^2>0$. Thus the tree does not expose an additional
public branch label that the original interface did not provide.
There are finitely many shapes and allocation labellings, although
their number is large. Write $a$ for the vector of their probabilities
and transfers and $\mathcal T(a)$ for these elementary constraints.

The following recursion produces the actual constraints and revenue
polynomials; it is not an optimization over an unspecified mechanism
space. At a node $h$ at date $t$, full actual history
$\eta=\theta^{1:t}$, and current report $r$, let
$F_{i,h,r,\eta}(a,z)$ be current-plus-future truthful utility, discounted
from date $t$. If edge $e$ after $r$ has data $(w_e,x_e,p_e)$ and child
$h_e$, then
\[
 F_{i,h,r,\eta}=
 \sum_e w_e\left[
 v_i(x_e,\eta_{it})-p_{ei}
 +\beta\sum_{\xi\in\Theta}
 K_{z,t}(\xi\mid\eta,x^{1:t-1},x_e)
 F_{i,h_e,\xi,\eta\xi}
 \right].
 \tag{F3}
\]
The last sum is omitted at $t=T$. Define $G_{h,r,\eta}$ by the same
recursion with current payoff $\sum_i p_{ei}$ in place of
$v_i(x_e,\eta_{it})-p_{ei}$. The root revenue polynomial is
\[
 R(a,z)=\sum_{\eta\in\Theta}\pi_z(\eta)\,
 G_{\varnothing,\eta,\eta}(a,z).
 \tag{F4}
\]

For agent $i$, a public node $h$, and an own private history
$\eta_i^{1:t}$, reconstruct the full actual past $\kappa$ by combining
its actual past with the opponents' past public reports. Write
$c=\eta_{it}$. The unnormalized probability of the current opponents'
type $\xi_{-i}$ is
\[
 \ell_z(\xi_{-i})=
 \begin{cases}
 \pi_z(c,\xi_{-i}),&t=1,\\
 K_{z,t-1}((c,\xi_{-i})\mid\kappa,x^{1:t-1}),&t>1.
 \end{cases}
 \tag{F5}
\]
Its sum is strictly positive. For every alternative current report
$c'\in\Theta_i$, form the polynomial
\[
 I_{i,h,\eta_i,c'}(a,z)=
 \sum_{\xi_{-i}}\ell_z(\xi_{-i})
 \left[
 F_{i,h,(c,\xi_{-i}),\,\kappa(c,\xi_{-i})}
 -
 F_{i,h,(c',\xi_{-i}),\,\kappa(c,\xi_{-i})}
 \right],
 \tag{F6}
\]
and the participation polynomial
\[
 J_{i,h,\eta_i}(a,z)=
 \sum_{\xi_{-i}}\ell_z(\xi_{-i})
 F_{i,h,(c,\xi_{-i}),\,\kappa(c,\xi_{-i})}.
 \tag{F7}
\]
At date one $\kappa$ is empty. All indices in these displays are finite.
Division by the positive conditional-probability denominator has been
removed, so no rational-function sign issue occurs.

\begin{proposition}[Exact finite feasibility test]\label{prop:finite-test}
A parameter vector $a$ defines a robustly truthful, voluntarily continued
mechanism precisely when $\mathcal T(a)$ holds and
\[
 \forall z\in A:\quad I_{i,h,\eta_i,c'}(a,z)\geq0,\qquad
 J_{i,h,\eta_i}(a,z)\geq0
 \quad\text{for every displayed finite index.}
 \tag{F8}
\]
For every law $z$, its revenue is \textup{(F4)}. Conversely every mechanism
in the original continuous-transfer class has the revenue function of
one of these finite feasible trees.
\end{proposition}
\begin{proof}
The conditional-probability denominator in \textup{(F5)} is positive,
so \textup{(F6)--(F7)} are exactly its one-step incentive and continuation
participation inequalities. The finite-horizon induction in
Lemma~\ref{lem:compression} proves equivalence with all adapted strategies.
The recursion proves the revenue assertion, and that lemma proves the
converse for the full original mechanism class.
\end{proof}

\section{Exact guarantees, positivity, and implementable formats}
\label{sec:elimination}

For a fixed tree shape write $\mathrm{Feas}(a;A)$ for the finite
first-order formula \textup{(F8)} together with $\mathcal T(a)$.
Disjunction over the finitely many shapes is understood in existential
statements, and conjunction over their separate parameter vectors in
universal statements. Thus no infinite-dimensional parameter or
unbounded discrete shape variable is hidden in a displayed quantifier.
The singleton formula $\mathrm{Feas}(a;\{z\})$ is obtained by evaluating
the same incentive and participation polynomials at that law.

There is no general compactness inference from bounded transfers:
the distinct-observable-outcome restrictions are strict inequalities.
Accordingly we encode a supremum, and separately decide whether it is
attained. The formula
\begin{align*}
 \mathrm{Bench}(z,b)\ \Longleftrightarrow\
 &\bigl[\forall a:\ 
       \mathrm{Feas}(a;\{z\})\ \Longrightarrow\ R(a,z)\leq b\bigr]\\
 {}\land&
 \bigl[\forall\varepsilon>0\ \exists a:\ 
       \mathrm{Feas}(a;\{z\})\ \land\ R(a,z)>b-\varepsilon\bigr]
 \tag{F9}
\end{align*}
is an explicit first-order formula defining $b=\mathrm{OPT}_z$.
The benchmark is finite because payments are bounded. Feasible
report-independent zero-transfer mechanisms ensure that its feasible
set is nonempty and its value is nonnegative.

For the ratio problem require $A\ne\varnothing$ and
$\inf_{z\in A}\mathrm{OPT}_z>0$, as in the maintained target.
The complete maintained-domain test is
\[
 [\exists z\,(z\in A)]\ \land\
 [\exists\delta>0\ \forall z,b:\quad
   [z\in A\land\mathrm{Bench}(z,b)]\Longrightarrow b\geq\delta].
 \tag{F10}
\]
If it fails, one reports that the maintained uniform-positive-benchmark
class fails, rather than dividing by zero. Individual zero-benchmark
laws are identified by $\mathrm{Bench}(z,0)$; removing them gives the
separately declared positive-benchmark ratio problem if desired.

For a real number $r$ define the finite algebraic feasibility formula
\begin{align*}
 \mathrm{Guarantee}(r)\ \Longleftrightarrow\ \exists a:\ 
 &\mathrm{Feas}(a;A)\\
 {}\land&
 \forall z,b:\quad
 [z\in A\land\mathrm{Bench}(z,b)]
 \Longrightarrow R(a,z)\geq rb .
 \tag{F11}
\end{align*}
Finally define
\begin{align*}
 \mathrm{Ratio}(u)\ \Longleftrightarrow\
 &[\forall r:\ \mathrm{Guarantee}(r)\Longrightarrow r\leq u]\\
 {}\land&
 [\forall\varepsilon>0\ \exists r:\ 
   \mathrm{Guarantee}(r)\land r>u-\varepsilon].
 \tag{F12}
\end{align*}
All quantifiers in \textup{(F8)--(F12)} range over finitely many real
coordinates or the real line. When $A$ has the specified semialgebraic
presentation, they can be eliminated by the quantifier-elimination theorem
for real closed fields~\cite{Collins}. Empty $A$ is reported separately by
\textup{(F10)}, rather than assigned a ratio.

\begin{theorem}[Complete finite algebraic characterization]
\label{thm:algebraic}
For every finite primitive model and every semialgebraically specified
ambiguity class satisfying the declared information and continuation
conventions, \textup{(F2)--(F12)} give the following.
\begin{enumerate}
\item The revenue-function set of all robust mechanisms, including all
continuous-transfer signals and all controlled transitions, is represented
exactly by the finite public-outcome format \textup{(F3)}.
\item The known-law benchmark and the robust ratio have exact first-order
real formulas. Under \textup{(F10)}, \textup{(F12)} has exactly one solution,
and that solution is the original $r^*(A)$.
\item Elimination produces a quantifier-free algebraic criterion, in the
input primitives, for a positive ratio, every proposed guarantee, and
attainment of the supremal ratio. If the exact optimum is not attained,
every strictly smaller ratio has a finite implementing tree. In particular,
for every $\varepsilon>0$ such a tree achieves a ratio exceeding
$r^*(A)-\varepsilon$.
\item Any declared finite-tree implementable format described by additional
finite polynomial constraints on these parameters can be treated in the
same way, after testing its feasible set for nonemptiness. Such formats
can impose deterministic outcomes, report-independent allocations,
history-independent per-date rules, or specified finite posted-price
or auction rules. Comparing their formulas decides exactly when that
declared format loses revenue or when independently heterogeneous
admissible beliefs remove a proposed intertemporal-screening gain.
To identify a larger restricted class of arbitrary continuous-signal
mechanisms with such a format would require a separate proof that
compression preserves that restriction; it is not assumed here.
\end{enumerate}
For rational or real-algebraic input data and rational or real-algebraic
requested guarantees or positive accuracies, these are terminating exact
algorithms with witness mechanisms whenever the corresponding
existential formula is feasible. No algorithm for an unencoded arbitrary
real input, and no polynomial-time bound, is asserted.
\end{theorem}
\begin{proof}
Proposition~\ref{prop:finite-test} proves the exact representation.
Its singleton version proves that \textup{(F9)} defines the original
known-law supremum, rather than a smaller mechanism class. A feasible
robust mechanism is feasible under each singleton law, so its revenue
ratio is at most one; a zero-transfer rule guarantees zero. Thus the
set defined by \textup{(F11)} is nonempty, bounded above, and downward
closed. Its supremum exists. By the exact representation,
\[
 \sup\{r:\mathrm{Guarantee}(r)\}
 =\sup_M\inf_{z\in A}\frac{\mathrm{Rev}_z(M)}{\mathrm{OPT}_z}.
\]
The upper-bound and approximation clauses in \textup{(F12)} characterize
this unique supremum without assuming that it is attained.

Real-closed-field quantifier elimination applies because every displayed
expression is polynomial and every discrete shape choice was enumerated
finitely. Applying it to $\mathrm{Ratio}(u)\land u>0$ gives the positivity
criterion; applying it to
$\mathrm{Ratio}(u)\land\mathrm{Guarantee}(u)$ gives the attainment test.
For $r<u$, the definition of a supremum gives an attainable guarantee
larger than $r$, hence $\mathrm{Guarantee}(r)$ by downward closure.
Over algebraic inputs, elimination and real-algebraic sample-point
selection return actual parameter vectors for the existential formulas.

For a format restriction, append its constraints before the existential
quantifier in \textup{(F11)} and use the same benchmark, if comparison is
against the original known-law optimum. Alternatively append them also
in \textup{(F9)} if that is the declared restricted benchmark. The distinction
is explicit. First test $\exists a\,\mathrm{Feas}(a;A)$ with those added
constraints, and report an empty restricted feasible set separately.
If it is nonempty, bounded transfers and the uniform positive benchmark
give a finite lower bound on every ratio, so the supremum formula remains
well-defined even if that restricted format has no zero-transfer rule.
Requiring incentive compatibility for every independently chosen belief
profile $(z_1,\ldots,z_n)\in A^n$ gives exactly the same robust
incentive inequalities: each agent's inequalities already hold for every
$z\in A$ separately. Thus the robust feasible set is unchanged when
robustness also permits these heterogeneous belief profiles. Comparing the unrestricted robust, restricted robust, or
singleton-law revenue formulas supplies finite algebraic tests for the
claimed screening gains. These are criteria in finite primitive
coordinates, obtained by an explicit elimination procedure, rather than
an unevaluated optimization over mechanisms.
\end{proof}

\paragraph{How the named ambiguity restrictions enter.}
For exogenous type paths, known one-date marginals and finite-path
moments are linear in joint path probabilities, and therefore polynomial
in the controlled-kernel coordinates of \textup{(F1)}. The recursions
and support bounds continue to use those kernel coordinates. With controlled transitions, finite-horizon path
probabilities are products of coordinates in \textup{(F1)}, so moments
and known marginals under declared fixed allocation interventions are
polynomial equations. Bounds on each transition coordinate are polynomial
inequalities. Any finite Boolean combination of such restrictions is
covered. A marginal defined only after the mechanism is chosen would
change the ambiguity class with the decision; it must instead be declared
under fixed interventions or encoded as an explicitly different model.
The algorithm keeps arbitrary correlations, type heterogeneity,
allocation-dependent transitions, and the specified feasible prefixes.

\paragraph{Distribution-free variant.}
For exogenous types, replace the Bayesian sums in \textup{(F6)--(F7)}
by the corresponding inequalities for every complete true type path,
with the future types fixed pathwise. With controlled types the
appropriate pathwise object is a complete nature response tree: for
every true-type and allocation prefix, it specifies the next type.
There are finitely many such response trees. A deviation then follows
its own allocation branch of that same response tree; one does not
hold an allocation-affected future type path artificially fixed.
Evaluate the continuation recursions at these zero--one controlled
kernels and impose the one-step and participation inequalities for
every compatible actual history and every such future response tree.
They remain finitely many polynomial inequalities, and the coefficient
preservation in Lemma~\ref{lem:compression} preserves all their values.
The same finite-horizon induction gives pathwise continuation IC.
Use these stronger constraints both for the candidate and the
known-law benchmark when that is the stated variant. This changes the
incentive feasible set; it does not silently identify ex-post and
Bayesian IC.


% Body component for the D82-46 manuscript; not a standalone solution claim.
\section{A distribution-independent positive guarantee}
\label{sec:positive-guarantee}

The following restrictions permit a strictly positive guarantee even when
allocations control future types. They are sufficient conditions, not
conditions imposed silently on the general finite optimization theorem.
Assume $n,T\geq1$, $0<\beta\leq1$, $0\leq v_i(x,\theta_i)\leq1$, and
payments in $[-B,B]$ with $B>0$. There is a null allocation $0\in X$ with
$v_i(0,\theta_i)=0$ for every agent and type. Let
$\mathcal S\subseteq X^T$ be the nonempty finite set of feasible allocation
paths. It is independent of types and satisfies
\[
 (x_1,\ldots,x_T)\in\mathcal S
 \ \Longrightarrow\
 (x_1,\ldots,x_k,0,\ldots,0)\in\mathcal S
 \quad(0\leq k\leq T).
 \tag{PG1}
\]
Thus termination after any feasible prefix is feasible. Write
$L=|\mathcal S|$. Initial laws and full-history controlled type-transition
kernels may be arbitrary within the maintained full-support model. In
particular their numerical probabilities need not be known to the designer.
They depend on actual types and realized allocations, as in the main model,
and are unaffected by reports once those quantities are fixed.

Refusal is available before reporting in each period, yields zero future
utility and leaves past transfers sunk. The designer commits to terminate
all future allocations and transfers following any exit. This is a feasible
allocation rule by \textup{(PG1)} and gives the exiting agent its specified
outside option. Exit notifications are public. Current allocation and payment
are executed together; exit
does not permit retention of a purchased allocation without its payment.
Opponents report truthfully and stay, as in the definition of robust IC and
revenue. Truthful continuation is voluntarily optimal, so this specifies an
admissible truthful participation selection even at indifference histories.

Let $v_{\min}$ be the minimum strictly positive number among
$v_i(x_t,\theta_i)$ with $(x_1,\ldots,x_T)\in\mathcal S$, $1\leq t\leq T$,
and all agents and types. If this set is nonempty, put
\[
 p=\tfrac12\min\{B,v_{\min}\}>0.
 \tag{PG2}
\]

\begin{proposition}\label{prop:positive-guarantee}
Under these restrictions there is one mechanism $R$, independent of the
law $P$, robustly truthful and voluntarily continued at every private
history, such that
\[
 \mathrm{Rev}_P(R)\ \geq\ \frac{p}{nTL}\,\mathrm{OPT}_P
 \qquad\text{for every admissible }P.
 \tag{PG3}
\]
Hence every ambiguity set with positive benchmarks, using the declared
convention for zero benchmarks, has
$r^*(\mathcal P)\geq\min\{B,v_{\min}\}/(2nTL)>0$.
This bound makes no claim to sharpness or to a size-independent constant.
If no positive value defining $v_{\min}$ exists, every benchmark is zero.
\end{proposition}

\begin{proof}
\emph{Construction.}
Privately draw a triple using randomization independent of initial types,
nature's transition randomization and agents' private coins:
$(\tau,j,s)$ uniformly from
$\{1,\ldots,T\}\times\{1,\ldots,n\}\times\mathcal S$.
Its law and the committed protocol are public; its realization is not
announced. Before date $\tau$, execute the whole allocation $s_t$ for free,
ignoring every report. At date $\tau$, execute the whole allocation
$s_\tau$ and charge only agent $j$ the price $p$ if
$v_j(s_\tau,r_j)\geq p$. Otherwise execute allocation $0$ and charge
nothing. After date $\tau$, execute allocation $0$ and charge nothing,
regardless of reports. Terminate immediately upon any exit.
Notice that the offered object is the whole feasible allocation, with all
other agents' shares free; no downward-closed coordinate projection of $X$
is required. Every realized path is a feasible prefix followed by zeros,
and all payments respect their bounds.

\emph{The hidden draw does not enlarge the information structure.}
This description uses only internal mechanism randomness. It is also
implementable as committed public-history stochastic kernels. For any
public report/outcome history let $C(h)$ be the finite set of triples
consistent with the allocations and payments already observed when fed
those reports. Conditional on $h$, the next kernel assigns an outcome the
fraction of triples in $C(h)$ that produce that outcome for the current
reports; afterwards update the consistency set. When $C(h)$ is empty use
permanent null continuation. No triple label or future random seed is
revealed. Given the actual type and allocation history, nature's likelihood
is the same for every compatible triple. An adapted reporting strategy
also uses only the agent's types, observed history and private coins, not
the unobserved triple. These likelihood factors cancel when conditioning,
so the described uniform-consistency kernels implement the private draw
under every admissible law and every such strategy. Equivalently, their
finite product along any fixed observable report/outcome route equals the
fraction of triples producing that route.

\emph{Incentives and continuation participation.}
Fix any private history, including one reached after earlier own lies,
and any compatible hidden triple. If the mechanism has already terminated
or its active date has passed, every
future allocation and payment is zero. If its active date is still ahead,
all preceding allocations are free, nonnegative in value and independent
of every report. Reports at those dates cannot change allocations or the
controlled type-transition law. If the active target is a different agent,
the considered agent's reports have no effect at any remaining date.
If this agent is the active target, its report at that date chooses between
utilities $v_j(s_\tau,\theta_{j\tau})-p$ and $0$, with no subsequent
utility. Truthful reporting chooses a maximizing alternative, for every
possible offered allocation. This holds whether the agent can infer the
active triple from its history or not. A lie at the active date may change
future types through the allocation, but those types have zero subsequent
utility because all later allocations and transfers are null.

Thus, fixing the hidden triple and coupling nature identically until the
first allocation-changing decision, truth dominates every adapted reporting
strategy. Exiting earlier gives zero continuation; truthful continuation is
nonnegative in every realization. Randomizing or learning from observed
outcomes cannot improve on these pointwise comparisons. Averaging under
any posterior over compatible triples and any law $P$ proves robust
continuation IC and participation at every private history. Inconsistent
public histories have null continuation, which has the same properties.
Sunk payments create no refund deviation.

\emph{A finite-prefix welfare bound.}
Fix $P$. For a type history
$\eta_t=(\theta^1,\ldots,\theta^t)$ and allocation prefix
$z_{1:t-1}$ write its nature factor as
\[
 d_P(\eta_t;z_{1:t-1})
 =q_1(\theta^1)\prod_{a=1}^{t-1}
 q_{a+1}(\theta^{a+1}\mid\theta^1,\ldots,\theta^a,z_{1:a}).
 \tag{PG4}
\]
Let $\operatorname{Pref}_t(\mathcal S)$ be the distinct feasible length-$t$
allocation prefixes. Define the finite, possibly loose upper bound
\[
 H_P=\sum_{t=1}^T\sum_{i=1}^n\beta^{t-1}
 \sum_{z_{1:t}\in\operatorname{Pref}_t(\mathcal S)}
 \sum_{\eta_t}
 d_P(\eta_t;z_{1:t-1})\,v_i(z_t,\theta_i^t).
 \tag{PG5}
\]
For any feasible known-law mechanism $M$, its expected discounted allocated
value is at most $H_P$. Indeed, conditional on a fixed full actual type
history, expand over the finitely many allocation prefixes and integrate
all payments and internal randomization. Each resulting joint route
probability is the factor \textup{(PG4)} times a nonnegative mechanism
likelihood at most one. Nature's factor can be pulled out of this integral
because its transitions depend on the fixed types and allocations, not on
hidden coins or payments. Nonnegativity of valuations therefore permits
replacing every mechanism likelihood by one, giving \textup{(PG5)}.
This argument allows arbitrarily adaptive payments and lotteries in $M$.

Initial voluntary participation bounds its revenue by its expected
allocated value: the sum of agents' expected utilities equals allocated
value minus transfers and is nonnegative. Consequently
$\mathrm{Rev}_P(M)\leq H_P$, and taking the supremum gives
$\mathrm{OPT}_P\leq H_P$. Attainment is unnecessary.

\emph{Revenue of the constructed mechanism.}
Before its active date the sampled allocation path is executed open loop.
Because every positive value is at least $v_{\min}\geq2p$, truthful
acceptance at date $t$ is equivalent to
$v_i(s_t,\theta_i^t)>0$. Hence
\[
 \mathrm{Rev}_P(R)
 =\frac{p}{nTL}
 \sum_{t=1}^T\sum_{i=1}^n\beta^{t-1}
 \sum_{s\in\mathcal S}\sum_{\eta_t}
 d_P(\eta_t;s_{1:t-1})
 \mathbf1\{v_i(s_t,\theta_i^t)>0\}.
 \tag{PG6}
\]
Every feasible length-$t$ prefix has at least one full extension in
$\mathcal S$. Moreover $\mathbf1\{v>0\}\geq v$ for $0\leq v\leq1$.
Thus the sum in \textup{(PG6)} is at least $H_P$, proving
$\mathrm{Rev}_P(R)\geq pH_P/(nTL)\geq p\mathrm{OPT}_P/(nTL)$.
This is \textup{(PG3)}. Full support implies strictly positive revenue
under each $P$ whenever a positive value occurs on a feasible path.
If all such values are zero, initial participation and the feasible null
mechanism instead give $\mathrm{OPT}_P=0$ for every $P$.
\end{proof}

The factor depends on the finite feasible path set and the smallest
positive value. Allowing these primitives to degenerate can make the bound
arbitrarily small. The proposition therefore identifies a concrete positive
domain; it does not imply a uniform positive bound when valuation spaces,
horizons, payment bounds or resource systems range without restriction.


\section{Known marginals: an exact loss from intertemporal screening}
\label{sec:intertemporal-example}

Consider one buyer and two dates, with discount factor one and a fresh
unit at each date. Its value is either $a=1/2$ or $b=1$, realized allocations lie in
$\{0,1\}$ (their expected probabilities lie in $[0,1]$), and payments satisfy $|p_t|\leq B$ for some $B\geq5/4$. The
buyer observes only its own current and past values and the public history.
Type evolution is exogenous. Exit is available before reporting at each
date, leaves past payments sunk, and gives zero future utility. The
ambiguity set consists of the full-support laws
\[
 P_\rho(a,a)=P_\rho(b,b)=\frac{1+\rho}{4},\qquad
 P_\rho(a,b)=P_\rho(b,a)=\frac{1-\rho}{4},\qquad -1<\rho<1.
 \tag{IT1}
\]
Both one-date marginals are known and uniform. The law $P_0$ is the
independent law.

\begin{proposition}\label{prop:intertemporal-loss}
Every mechanism that is dynamically truthful and voluntarily continued
under every law in \textup{(IT1)} earns at most $1$ under $P_0$. This bound
is attained by charging $1/2$ for each unit independently of the history.
Under the known law $P_0$, a dynamically truthful mechanism with voluntary
continuation earns $9/8$. Consequently the robust minimax ratio for this
ambiguity set is at most $8/9$. In particular, known marginals and full
support do not preserve all gains from intertemporal screening.
\end{proposition}

\begin{proof}
Fix an arbitrary robust mechanism, including arbitrary randomized
allocations, payments, and continuation rules. For a first-date report
$r\in\{a,b\}$, let $x_r$ and $p_r$ denote its expected first allocation
and payment. For a second-date true value $z\in\{a,b\}$, let $q_r(z)$
and $U_r(z)$ denote the expected second allocation and truthful utility,
averaging over all first-date public outcomes and subsequent mechanism
randomization. These expectations are defined by the report input $r$
and the second value $z$, rather than by the buyer's true first value.
This is valid because the mechanism observes reports and its own outcomes,
not that hidden value, and the type process is exogenous.

At every realized first-date public outcome, second-date truthfulness
gives the usual two-type incentive inequalities. Averaging the inequality
for true value $b$ against report $a$ gives
\[
 d_r:=U_r(b)-U_r(a)\ \geq\ (b-a)q_r(a)=\tfrac12q_r(a).
 \tag{IT2}
\]
Continuation participation gives $U_r(a),U_r(b)\geq0$. In particular
$U_r(b)\geq U_r(a)$.

For a fixed first true value, its conditional second-value distribution
under $P_\rho$ converges to each of the two point masses as $\rho$ tends
to $1$ and $-1$. First-date robust incentive and participation inequalities
pass to both limits: there are only two future values and all payoffs are
bounded. Truthfulness for first value $b$, compared with reporting $a$
and then truthfully reporting the second value, therefore implies
\[
 p_b-p_a\leq (x_b-x_a)+U_b(z)-U_a(z)
 \qquad(z=a,b).
 \tag{IT3}
\]
First-date participation for true value $a$, using the limit with second
value $a$, implies
\[
 p_a\leq\tfrac12x_a+U_a(a).
 \tag{IT4}
\]
Only admissible full-support laws are used before taking these limits;
no degenerate law is inserted into the ambiguity set.

Put $u_r=U_r(a)$ and $\Delta x=x_b-x_a$. Combining \textup{(IT3)--(IT4)}
gives
\[
 \frac{p_a+p_b}{2}
 \leq \frac{x_b}{2}+u_a+\frac12
 \min\{u_b-u_a,\ u_b-u_a+d_b-d_a\}.
 \tag{IT5}
\]
At $P_0$, both the first type and the second type are independent and
uniform. Second-date expected payments, by the definition of utility,
are
\[
 \frac18\{q_a(a)+q_b(a)\}
 +\frac14\{q_a(b)+q_b(b)\}
 -\frac14\{2u_a+2u_b+d_a+d_b\}.
 \tag{IT6}
\]
Adding \textup{(IT5)} and \textup{(IT6)} and cancelling $u_a,u_b$ yields
\begin{align*}
 \mathrm{Rev}_{P_0}(M)
 &\leq \frac{x_b}{2}
   +\frac18\{q_a(a)+q_b(a)\}
   +\frac14\{q_a(b)+q_b(b)\}
   -\frac14(d_a+d_b)-\frac12(d_a-d_b)_+\\
 &\leq \frac{x_b}{2}+\frac14\{q_a(b)+q_b(b)\}\leq1.
 \tag{IT7}
\end{align*}
The middle inequality is \textup{(IT2)}; the last uses allocation
probabilities at most one. This argument applies to every public-payment
signal and every history-dependent randomized continuation, because only
averaged valid incentive inequalities were used. Charging $1/2$ in each
period is pathwise truthful and voluntarily continued, and earns exactly
$1$ under every law in \textup{(IT1)}.

For the known independent law, allocate the first unit for either report.
Charge $1/2$ after report $a$, and offer the second unit at price $1$.
Charge $5/4$ after report $b$, and provide the second unit for free.
At date two both rules are truthful and yield nonnegative utility at
each realized value. At date one the independent future value has mean
$3/4$. True first value $a$ obtains utility zero from either first report:
the $a$ report gives $a-1/2=0$ and zero future rent, while the $b$ report
gives $a-5/4+3/4=0$. True first value $b$ obtains $1/2$ from either report.
Thus truthful reporting and continuation are optimal, including against
adapted reporting-and-exit strategies by backward induction. All payments
respect the stated bound. Expected revenue is
\[
 \frac12\left(\frac12+\frac54\right)
 +\Pr\{\theta_1=a,\theta_2=b\}\cdot1
 =\frac78+\frac14=\frac98.
 \tag{IT8}
\]
Therefore $\mathrm{OPT}_{P_0}\geq9/8$, while every robust mechanism earns
at most $1$ at $P_0$. The ratio bound follows. This proves an upper bound
on the exact robust ratio, not an assertion that $8/9$ is its exact value.
\end{proof}


\begin{thebibliography}{9}
\bibitem{BV}
Dirk Bergemann and Juuso V\"alim\"aki.
\emph{Dynamic Mechanism Design: An Introduction}.
Journal of Economic Literature 57(2) (2019), 235--274.
Inspected Cowles Discussion Paper 2102R, June 19, 2018, Section~8,
printed pp.~48--50 (physical PDF pp.~49--51).
\url{https://cowles.yale.edu/sites/default/files/2022-09/d2102-r.pdf}.
\url{https://doi.org/10.1257/jel.20180892}.
\bibitem{Collins}
George E. Collins.
\emph{Quantifier Elimination for Real Closed Fields by Cylindrical
Algebraic Decomposition}.
In H. Brakhage (ed.), \emph{Automata Theory and Formal Languages},
Lecture Notes in Computer Science 33 (Springer, 1975), 134--183.
Section~3, Theorem~8 p.~158; algorithms EVAL and ELIM, pp.~159--160.
\url{https://doi.org/10.1007/3-540-07407-4_17}.
\end{thebibliography}
\end{document}
